% P-I01 | motivation: phflo-marked scIntDesign143-145; scTab donor-diversity experiment
Single-cell atlases provide increasingly broad reference maps, yet annotating a new experiment often requires inference from a small labelled set collected under different biological and technical conditions. More profiled cells do not necessarily supply more independent donors or more varied experimental settings. In supervised single-cell annotation, increasing donor diversity can contribute more to generalization than increasing cell count alone\cite{fischer2024sctab}. This distinction motivates models that learn how to interpret limited experimental context across heterogeneous tasks.

% P-I02 | model successes and bounded critique; do not infer biology is linear
Single-cell foundation models such as Geneformer, scGPT, scFoundation, UCE and TranscriptFormer learn reusable representations from large collections of measured profiles\cite{theodoris2023geneformer,cui2024scgpt,hao2024scfoundation,rosen2026uce,pearce2026transcriptformer}. Their breadth supports transfer across biological settings, but large-scale pretraining does not guarantee an advantage on each downstream task. Evaluations have found limitations in frozen representations and saturation as pretraining collections expand\cite{kedzierska2025zeroshot,denadel2026scaling}; recent preprint evidence also shows that normalization and linear methods can remain competitive across several benchmarks\cite{souza2026parameterfree}. These findings motivate a complementary question: can learning improve the inference rule that combines a few labelled examples, even when the expression measurements themselves already provide useful features?

% P-I03 | PFN relation and closest cellular prior work; no first biologicalPFN claim
Meta-learning already supports transfer across heterogeneous single-cell experiments, as demonstrated by MARS\cite{brbic2020mars}. Prior-data fitted networks (PFNs) provide a complementary way to learn inference rules. They are trained across prediction problems drawn from a generative prior, amortizing support-conditioned inference into a forward pass\cite{muller2022pfn}. TabPFN and TabICL demonstrate that synthetic task distributions can produce effective general-purpose tabular learners\cite{hollmann2023tabpfn,hollmann2025tabpfn,qu2025tabicl,qu2026tabiclv2}. Recent MapPFN and PerturbPFN work further applies prior-based learning to cellular perturbation prediction\cite{sextro2026mappfn,gao2026perturbpfn}. For transcriptomic classification, the opportunity is to encode the variation encountered across cellular experiments directly into the training prior, while accommodating thousands of sparse features and scarce labels.

% P-I04 | recovered simulator motivation; fair acknowledgement of configurable predecessors
Single-cell simulators already model count distributions, dependence, donor variation and regulatory dynamics through complementary statistical and mechanistic approaches\cite{zappia2017splatter,zhang2019symsim,sun2021scdesign2,song2024scdesign3,azodi2021splatpop,dibaeinia2020sergio,cannoodt2021dyngen,li2025scmultisim}. Learned generative models also support conditional synthesis and augmentation\cite{marouf2020cscgan,luo2024scdiffusion}. For PFN training, the design objective extends beyond reproducing a reference dataset to generating diverse labelled inference problems. Controlled variation in signal strength, hidden subpopulations, donor effects and measurement noise can expose a learner to experimental configurations without requiring each configuration to be measured and annotated. Synthetic-to-real learning in optical flow and face analysis provides precedents for learning from known simulated structure\cite{dosovitskiy2015flownet,wood2021fake}. Nevertheless, simulator fidelity depends on the statistic and scenario being evaluated\cite{crowell2023shaky,duo2024systematic}. The relevant test is whether the induced prior teaches inference that transfers to real biological data.

% P-I05 | this study: two uses of one prior, support-conditioned gene anonymity
Here we introduce CellPFN and its procedural cellular prior. CellPFN learns from synthetic tasks while preserving gene correspondence within each task and omitting persistent gene identities across tasks. We test two consequences of this prior: whether a compact cellular model can achieve competitive biological inference after substantially fewer training episodes than a general tabular model, and whether tasks from the same simulator can improve released tabular models through continued training. We further examine transcriptome-scale inputs and efficient query-gene explanations. The aim is a reusable method for learning from limited labelled biological context, with the prior's usefulness established through downstream inference rather than simulation scale alone.
